\documentclass[11pt]{article}
\usepackage[T1]{fontenc}
\usepackage{textcomp}
\usepackage[margin=0.75in]{geometry}
\usepackage{amsmath,amssymb,amsthm}
\usepackage{array,booktabs}
\usepackage{hyperref}
\setlength{\parindent}{0pt}
\newtheorem{theorem}{Theorem}
\theoremstyle{definition}
\newtheorem{definition}{Definition}
\title{Flow Proof Artifact: prop-32-complements-and-applied-area}
\author{Generated by \texttt{flow doc proof}}
\date{Flow Proof Book}
\begin{document}
\maketitle
If ABCD is a parallelogram and BE is a parallelogram about the diameter, the complement CD is such that the whole is to the complement as the other about the diameter is to CD.\\[0.5em]
\textbf{Source.} euclid — Elements, Book VI, Proposition 32\\[0.5em]
\bigskip
\noindent{\large \textbf{Derived fact 1.} \textit{If ABCD is a parallelogram and BE is a parallelogram about the diameter, the complement CD is such that the whole is to the complement as the other about the diameter is to CD} \hfill the Euclidean plane \cdot Euclid Book VI \cdot Proposition 32: complements with an applied parallelogram relate by duplicate ratio}\par
\smallskip\noindent\textit{Coordinate: the Euclidean plane \cdot Euclid Book VI \cdot Proposition 32: complements with an applied parallelogram relate by duplicate ratio}\par
\smallskip\noindent\textit{Source: euclid — Elements, Book VI, Proposition 32}\par
\smallskip\noindent\textit{Built on: proposition 24: parallelograms about a diameter are similar to the whole and to one another, for Euclid Book VI on the Euclidean plane, proposition 31: similar parallelograms are as the duplicate ratio of homologous sides, for Euclid Book VI on the Euclidean plane}\par
\medskip
\noindent\textbf{Goal.} If ABCD is a parallelogram and BE is a parallelogram about the diameter, the complement CD is such that the whole is to the complement as the other about the diameter is to CD\par
\noindent$\displaystyle \text{complement applied area ratio}$\par
\medskip
\noindent
\begin{tabular}{@{} >{\raggedright\arraybackslash}p{0.06\textwidth} >{\raggedright\arraybackslash}p{0.47\textwidth} >{\raggedleft\arraybackslash}p{0.41\textwidth} @{}}
\textbf{\#} & \textbf{Proof} & \textbf{Mathematics} \\
\midrule
\textbf{1.} & We prove that proposition 32: complements with an applied parallelogram relate by duplicate ratio for Euclid Book VI the Euclidean plane. & \\[0.45em]
\textbf{2.} & Let ABCD = whole\_parallelogram. & \\[0.45em]
\textbf{3.} & Let BE = about\_diameter. & \\[0.45em]
\textbf{4.} & Let CD = complement. & \\[0.45em]
\textbf{5.} & From step 2, step 3, and step 4, we can deduce that ratio whole complement holds. & $\displaystyle \text{ratio whole complement holds}$ \\[0.45em]
\textbf{6.} & We invoke the derived fact governing Euclid Book VI the Euclidean plane: proposition 24: parallelograms about a diameter are similar to the whole and to one another, for Euclid Book VI on the Euclidean plane. & \\[0.45em]
\textbf{7.} & We invoke the derived fact governing Euclid Book VI the Euclidean plane: proposition 31: similar parallelograms are as the duplicate ratio of homologous sides, for Euclid Book VI on the Euclidean plane. & \\[0.45em]
\textbf{8.} & From step 6 and step 7, this implies complement applied area ratio. Hence proven. & $\displaystyle \text{complement applied area ratio}$ \\[0.45em]
\bottomrule
\end{tabular}
\label{thm:Geometry-Euclid Book VI-proposition_32_complements_with_an_applied_parallelogram_relate_by_duplicate_ratio}
\par\medskip\hrule
\end{document}
